\clearpage
\fancyhead[L]{\small IRIS.}
\SourceLeaf{478}
\begin{center}
\rule{0.55\linewidth}{0.4pt}\\[2pt]
Nachrichten der Königl. Gesellsch. d. Wiss. zu Göttingen. No. 11. 13 September 1847.\\[2pt]
\rule{0.55\linewidth}{0.4pt}
\end{center}

I first learned from a Paris newspaper dated 16 August that Mr \textsc{Hind} had discovered a new planet in London on 13 August. On the strength of that report I observed it on the meridian on 21 August. Since the report had given only two positions an hour apart, there was at first some uncertainty about its identity. This was soon dispelled when further English observations from 14 and 15 August arrived. Since then I have observed the planet, which, as is well known, has been named Iris, with the meridian circle on every evening the weather has allowed. I have obtained the following positions so far.

\begin{center}
\setlength{\tabcolsep}{4pt}
\begin{tabular}{lr|r|r|r}
\multicolumn{2}{c|}{1847} & \shortstack{Mean time\\in Göttingen} & \shortstack{Right\\ascension} & \shortstack{Southern\\declination} \\\hline
August & 21 & $9^{\rm u}\ 52^{\rm m}\ 54^{\rm s}1$ & $297^\circ\ 47'\ 20''8$ & $13^\circ\ 42'\ 19''9$ \\
 & 28 & $9\ 21\ 16.6$ & $296\ 45\ 38.2$ & $13\ 54\ 47.9$ \\
 & 29 & $9\ 16\ 53.3$ & $296\ 38\ 46.6$ & $13\ 56\ 34.2$ \\
Sept. & 4 & $8\ 51\ 12.0$ & $296\ 7\ 14.1$ & $14\ 6\ 3.8$ \\
 & 6 & $8\ 42\ 54.5$ & $296\ 0\ 46.5$ & $14\ 8\ 56.1$ \\
\end{tabular}
\end{center}

In these observations (with atmospheric conditions generally less than favourable), the planet appeared only as bright as a star of the 9th magnitude.

As soon as the observation of 28 August had been reduced, Professor \textsc{Goldschmidt} set about calculating the elliptic elements. He based them on my two observations of 21 and 28 August and one made in Cambridge by Mr \textsc{Challis} on 14 August. The following are his results.

\begin{center}
\small
\setlength{\tabcolsep}{2.5pt}
\begin{tabular}{llr}
\multicolumn{2}{l}{\shortstack[l]{Mean anomaly 1847 August 14. $0^{\rm u}$,\\Mean Göttingen time}} & $288^\circ\ 53'\ 21''9$ \\
Longitude of perihelion & \mbox{.\ .\ .\ .\ .\ .\ .\ .\ .\ .\ .\ .\ .\ .} & $44\ 6\ 49.3$ \\
Longitude of the ascending node & \mbox{.\ .\ .\ .\ .\ .\ .\ .\ .} & $260\ 27\ 58.8$ \\
\multicolumn{3}{l}{both reckoned from the apparent equinox of 14 August} \\
Mean daily sidereal motion & \mbox{.\ .\ .\ .\ .\ .\ .} & $982''5907$ \\
Logarithm of the semimajor axis & \mbox{.\ .\ .\ .\ .\ .\ .\ .} & $0.3717560$ \\
Eccentricity & \mbox{.\ .\ .\ .\ .\ .\ .\ .\ .\ .\ .\ .\ .\ .\ .\ .} & $0.2135543$ \\
Inclination of the orbit & \mbox{.\ .\ .\ .\ .\ .\ .\ .\ .\ .\ .\ .\ .\ .} & $5^\circ\ 29'\ 55''8$ \\
\end{tabular}
\end{center}

Mr \textsc{Gould} of Boston, who is pursuing his astronomical studies here with great zeal, has likewise calculated an elliptic orbit from the same observations. It differs little from the one given above and also still agrees very well with the later observations made so far.

Two conclusions can already be drawn from this. First, it is now certain that Iris is a planet and that its orbit, like those of Ceres, Pallas, Juno, Vesta, Astraea and Hebe, lies between the orbits of Mars and Jupiter. Second, the planet will probably be considerably brighter at its apparition in 1848 than it is this year.

\noindent Göttingen, 7 September 1847.\hfill\textsc{Gauss}.
